\section{Architecture}
\label{sec:architecture}

To guarantee the globally unique mapping of human-readable names to machine identifiers, a naming system requires an agreed-upon history of requests, a strict set of rules, and a final registry of owners.

\zns{} treats this final registry as derived state by using Zcash exclusively for the history of requests, executing rules off-chain, and proving their enforcement on-chain.

\subsection{System Components}

The protocol relies on four distinct components:

\begin{description}
  \item[The Mint.] An off-chain execution environment. The Mint evaluates rules against user requests. It runs entirely inside a Trusted Execution Environment (TEE). It bundles a local Zcash node to read chain state trustlessly.
  \item[The Name Note.] A cryptographic proof of state transition. When the Mint evaluates a valid input, it publishes a Name Note to the Zcash blockchain. The Name Note binds a specific name to a specific owner.
  \item[The Attestation.] A hardware-generated proof of integrity. The TEE signs a remote attestation proving that the Mint's signing key is exclusively controlled by the approved protocol code.
  \item[The Resolver.] A permissionless client. The Resolver reads the immutable sequence of Zcash transactions. It verifies the Name Notes and deterministically computes the final derived state.
\end{description}

\subsection{Attestation Constraints}

The Resolver accepts a Mint execution if and only if the attestation proves five facts:

\begin{enumerate}
  \item The enclave code measurement matches the published audit manifest.
  \item The enclave holds the authoritative Mint spending key.
  \item The execution enforces the correct \zns{} protocol version.
  \item The execution targets the correct Zcash network.
  \item The enclave implements replay prevention via a protected persistent-state mechanism.
\end{enumerate}

The Resolver discards the output if an attestation fails any constraint. The Resolver discards the output if the TEE debug flag is enabled.

The active manifest must identify the immutable source code revision for every approved measurement. The active manifest must identify the public audit report for every approved measurement.

The rules for activation, measurement-approval, rollback-protection, upgrades, and key-rotation are \openstatus{}. This document defines the intended policy construction. It does not yet define the byte-level attestation check.

\subsection{Independence of Facts}

The security of the derived state relies on the strict orthogonality of its proofs. To accept a state transition, the Resolver must independently verify six distinct facts:

\begin{itemize}
  \item \textbf{Zcash validation:} The transaction satisfies Zcash consensus.
  \item \textbf{Mint self-send:} The recognized Mint spending authority authorized the transaction.
  \item \textbf{TEE attestation:} The recognized Mint identity executed the approved enclave code.
  \item \textbf{Name Note verification:} One name binding tuple matches one Orchard note commitment.
  \item \textbf{User authorization:} The user authenticated their intent via the active authorization mechanism.
  \item \textbf{Resolver replay:} The derived state matches the canonical log under a fixed rule set.
\end{itemize}